\documentclass[12pt,a4paper]{article}

\usepackage[T2A]{fontenc}
\usepackage[utf8]{inputenc}
\usepackage[russian]{babel}
\usepackage{geometry}
\usepackage{hyperref}

\geometry{margin=2.5cm}
\hypersetup{colorlinks=true, urlcolor=blue}

\begin{document}

\begin{titlepage}

\begin{center}
    \LARGE Национальный исследовательский университет ИТМО\\[8mm]

    Факультет Программной Инженерии и Компьютерной техники\\[37mm]

    \huge \textbf{NoSQL}\\[2mm]
    \textbf{Лабораторная работа №1} \\[2mm]
\end{center}

\vfill

\begin{flushright}
    \large Выполнили:

    Кузьмин Дмитрий Анатольевич, \\
    Софьин Вячеслав Евгеньевич \\[3mm]
    Группа:

    P3309\\[3mm]

    Преподаватели:

    Деменев Тимур Гурбанович,\\
    Коновалов Арсений Антонович
\end{flushright}

\vfill

\begin{center}
    Санкт-Петербург 2026
\end{center}

\end{titlepage}

\section*{Задание}

Вариант №1865355

Разработать веб-приложение на языке Java с использованием Spring Boot, реализующее каталог товаров.

Предметная область: интернет-магазин.
Основной объект системы: товар.
Пользовательская роль: преподаватель.
Обязательный сценарий: просмотр каталога и добавление в избранное.

Система хранения типа «ключ — значение»: Riak KV.

\subsection*{Дополнительные параметры}

\begin{itemize}
    \item сценарий временного хранения: ссылка восстановления доступа;
    \item сценарий кэширования: история операций;
    \item атомарный механизм: счётчик созданных заявок;
    \item исследование поведения системы: восстановление постоянных данных.
\end{itemize}

\subsection*{Требования к программе}

\begin{enumerate}
    \item Развернуть систему хранения, указанную в варианте.
    \item Реализовать подключение к ней из Java/Spring Boot.
    \item Выбрать и обосновать способ представления данных в модели «ключ — значение».
    \item Реализовать хранение части данных приложения в выбранной системе.
    \item Реализовать механизм временных данных: TTL для Redis, application-level expiration для Riak KV или lease для Etcd.
    \item Реализовать кэширование часто запрашиваемых данных.
    \item Реализовать назначенный атомарный счётчик, ограничение запросов или блокировку.
    \item Настроить сохранение и восстановление данных, если выбранная система поддерживает этот механизм.
    \item Исследовать назначенный сценарием режим поведения системы.
    \item Сделать выводы о применимости выбранного хранилища для разработанного приложения.
\end{enumerate}

\section*{Репозиторий с приложением}

\url{https://github.com/Dmitry-Kuzmin456/NoSQL_lab1}

\section*{Выводы}

Каталог, пользователи и заказы хранятся в PostgreSQL: по ним нужны фильтры, сортировка и страницы. В Riak KV вынесена та часть данных, которую удобно брать по одному ключу.

Избранное и корзина записаны как CRDT-карты с ключом пользователя: поле одного товара меняется отдельно и не затирает остальные. Сессии и токен восстановления пароля --- обычные пары «ключ — значение». У токена нет TTL на стороне Riak, срок 15 минут проверяет приложение и удаляет запись после использования или по истечении времени. История операций постоянно лежит в PostgreSQL, а в Riak кэшируется первая страница (cache-aside): при новой операции кэш сбрасывается. Число созданных заявок считается CRDT-счётчиком: приложение посылает приращение, а не перезаписывает число целиком.

Постоянные данные Riak переживают перезапуск контейнера за счёт тома Docker. Отдельно каталог \texttt{/var/lib/riak} снимается в архив и разворачивается обратно скриптами \texttt{backup\_riak.sh} и \texttt{restore\_riak.sh}. После восстановления возвращаются избранное, корзина и счётчики. Сессии, токены восстановления и кэш истории можно не восстанавливать: кэш собирается из PostgreSQL заново.

Riak уместен для избранного, корзины, сессий, одноразовых ссылок и счётчика. Для каталога и заявок он не подходит: нет выборок по цене, статусу и сортировке. Срок жизни записей приходится контролировать в коде. Копия файлов одного узла восстанавливает данные, но не заменяет репликацию кластера и не даёт такой же надёжности, как транзакции PostgreSQL. Для этого приложения Riak --- дополнение к реляционной базе под сценарии варианта, а не единственное хранилище.

\end{document}
